% Auto-generated from book/08_class_group.md — do not edit by hand
% Regenerated by scripts/md_to_latex.py

\chapter{Chapter 8 --- Composition and Class Groups in the Quaternionic Setting}
\label{ch:ch08_class_group}

This chapter lifts Gauss composition of binary quadratic forms and the ideal class group (Hatcher Chapters 7--8) to the gauged Hopf lattice and quaternion orders. We define composition of flux configurations / flywheels, construct class-group analogues, and examine how these structures organize Magic Islands and equivalence classes. The theory remains a \textbf{Model} construction, with \textbf{Open Problem 6} as its central open edge.

\textbf{Learning goals}

\begin{enumerate}
\item Recall Gauss composition of quadratic forms and its arithmetic meaning.  
\item Lift composition to flux configurations and flywheels on the gauged Hopf lattice.  
\item Construct class-group analogues for flux topographs (and, later, quaternion orders).  
\item Connect class-group structure to Magic Island counts and equivalence classes.  
\item Prepare the ground for quaternion algebras and ideal theory (Chapter 9).
\end{enumerate}

\textbf{Figures in this chapter}

\begin{center}
\small
\begin{tabular}{@{}lp{0.28\textwidth}p{0.28\textwidth}@{}}
\toprule
Tag & File & Role \\
\midrule
Fig. 8.1 & \texttt{figures/fig8\_1\_gauss\_composition.png} & Classical Gauss composition diagram \\
Fig. 8.2 & \texttt{figures/fig8\_2\_flywheel\_composition.png} & Composition of two flux topographs \\
Fig. 8.3 & \texttt{figures/fig8\_3\_class\_group\_analogue.png} & Class-group analogue on the lattice \\
Fig. 8.4 & \texttt{figures/fig8\_4\_island\_from\_class.png} & Magic Island from class-group data \\
Aux A8.1 & \texttt{figures/aux8\_1\_composition\_table.png} & Small composition table \\
\bottomrule
\end{tabular}
\end{center}


\textbf{Claim discipline}

\begin{center}
\small
\begin{tabular}{@{}lp{0.55\textwidth}@{}}
\toprule
Claim & Type \\
\midrule
Classical Gauss composition and ideal class groups (Hatcher Ch. 7--8) & \textbf{Theorem} (cited) \\
Composition of flux configurations / flywheels; class-group analogues & \textbf{Model} (core of OP6) \\
\texttt{qga/lib/composition.py} helpers & \textbf{Software fact} \\
\bottomrule
\end{tabular}
\end{center}


\bigskip
\noindent\rule{\textwidth}{0.4pt}
\bigskip

\section{8.1 Gauss composition --- classical reminder}
\label{ch:ch08_class_group:8-1-gauss-composition-classical-reminder}

In Hatcher Chapter 7, Gauss defined a composition law on (primitive) binary quadratic forms of the same discriminant. Up to equivalence, the law is associative and turns the set of classes into a finite abelian group---the \textbf{class group} of the discriminant. The same group appears as the ideal class group of the corresponding quadratic order (Hatcher Chapter 8).

\begin{figure}[htbp]
  \centering
  \includegraphics[width=0.92\textwidth,height=0.42\textheight,keepaspectratio]{fig8_1_gauss_composition.png}
  \caption{Figure 8.1 --- Classical Gauss composition.}
  \label{fig:fig8-1-gauss-composition}
\end{figure}

\begin{quote}\small\textit{Figure 8.1.* Forms \(f_1,f_2\) of fixed discriminant compose to a third form; classes form a finite abelian group. This is classical \textbf{Theorem} territory; we cite Hatcher / Gauss rather than re-prove.}\end{quote}

Why it matters here: class groups encode failure of unique factorization and control representation theory---exactly the arithmetic depth we want to lift for Magic Islands and the \(Z\mapsto\) map.

\bigskip
\noindent\rule{\textwidth}{0.4pt}
\bigskip

\section{8.2 Lifting composition to flux configurations and flywheels}
\label{ch:ch08_class_group:8-2-lifting-composition-to-flux-configurations-and-flywheels}

On the gauged Hopf lattice we define a \textbf{composition} of two flux topographs (or flywheels) by combining values and/or points, then reducing with the Chapter 6 normalization.

\subsection{Candidate methods (\textbf{Model})}
\label{ch:ch08_class_group:candidate-methods-model}

\begin{center}
\small
\begin{tabular}{@{}lp{0.55\textwidth}@{}}
\toprule
Method (code) & Rule \\
\midrule
\texttt{value\_sum} & Pointwise average of values; averaged then renormed points \\
\texttt{value\_product} & Pointwise product (scaled) \\
\texttt{value\_bilinear} & Soft bilinear blend \((V_a+V_b+V_a V_b)/3\) \\
\texttt{point\_mult} & Quaternion multiply corresponding lattice points; recompute functional \\
\bottomrule
\end{tabular}
\end{center}


\begin{lstlisting}[style=qga]
compose_flywheels(topo_a, topo_b, method="value_sum")
reduce_composition(composed) -> topograph, classification, magic_island_score
\end{lstlisting}

\begin{figure}[htbp]
  \centering
  \includegraphics[width=0.92\textwidth,height=0.42\textheight,keepaspectratio]{fig8_2_flywheel_composition.png}
  \caption{Figure 8.2 --- Composition of two flux topographs.}
  \label{fig:fig8-2-flywheel-composition}
\end{figure}

\begin{quote}\small\textit{Figure 8.2.* \(\Phi_1\), \(\Phi_2\), and \(\mathrm{reduce}(\Phi_1*\Phi_2)\) for the pedagogical \texttt{value\_sum} method.}\end{quote}

\textbf{Honest status.} On current candidate adjacency and modest angle lattices, composition tables often show \textbf{poor closure} and \textbf{low associativity scores} (\textbf{Software fact} from the sandbox). That is not a failure of the chapter---it is the experimental content of \textbf{Open Problem 6}.

\textbf{Open Problem 6 (core of this chapter).} Does a well-defined, associative composition law exist on equivalence classes of flux configurations / flywheels that:

\begin{enumerate}
\item Reduces to classical Gauss composition under a suitable projection or restriction?  
\item Is compatible with the gauge actions of Chapter 4?  
\item Produces a finite (or finitely generated) abelian group whose order/structure predicts Magic Island counts?
\end{enumerate}

Sandbox: \texttt{qga/lib/composition.py}. Depends on OP1--OP3.

\bigskip
\noindent\rule{\textwidth}{0.4pt}
\bigskip

\section{8.3 Class-group analogues}
\label{ch:ch08_class_group:8-3-class-group-analogues}

The set of equivalence classes of reduced flux topographs, equipped with a composition law, forms a \textbf{class-group analogue}. Its order is a \textbf{class-number analogue}---related to Chapter 6's \texttt{class\_number\_analogue}, now enriched by composition tables and associativity diagnostics.

\begin{lstlisting}[style=qga]
class_group_analogue(topo_or_list, method=..., dedup_tol=...)
  -> order, structure, table, closure_fraction, associativity, representatives

composition_table(reduced_set)
is_associative_up_to_equivalence(reduced_set, samples=...)
\end{lstlisting}

\begin{figure}[htbp]
  \centering
  \includegraphics[width=0.92\textwidth,height=0.42\textheight,keepaspectratio]{fig8_3_class_group_analogue.png}
  \caption{Figure 8.3 --- Class-group analogue.}
  \label{fig:fig8-3-class-group-analogue}
\end{figure}

\begin{quote}\small\textit{Figure 8.3.* Equivalence classes with composition arrows, plus a numerical composition table. \texttt{structure} strings such as \texttt{partial\_magma} or \texttt{open\_or\_non\_closed} flag OP6 incompleteness.}\end{quote}

Magic Islands can be viewed as \textbf{distinguished classes} (or clusters of classes) where composition yields high \texttt{magic\_island\_score} or high portal \texttt{stability\_score}. Predicting islands from group structure alone is OP3 ? OP6.

\begin{figure}[htbp]
  \centering
  \includegraphics[width=0.92\textwidth,height=0.42\textheight,keepaspectratio]{fig8_4_island_from_class.png}
  \caption{Figure 8.4 --- Magic Island from class-group data.}
  \label{fig:fig8-4-island-from-class}
\end{figure}

\begin{quote}\small\textit{Figure 8.4.* Schematic: class-group analogue \(\rightarrow\) distinguished high-stability region. \textbf{Model} narrative, not a derived theorem.}\end{quote}

\begin{figure}[htbp]
  \centering
  \includegraphics[width=0.92\textwidth,height=0.42\textheight,keepaspectratio]{aux8_1_composition_table.png}
  \caption{Auxiliary Figure A8.1 --- Composition table.}
  \label{fig:aux8-1-composition-table}
\end{figure}

\begin{quote}\small\textit{Auxiliary Figure A8.1.* Small \(n\times n\) table of reduced class indices under composition. Entries \(-1\) mean ``not closed in the set within tolerance.''}\end{quote}

\bigskip
\noindent\rule{\textwidth}{0.4pt}
\bigskip

\section{8.4 First computational labs}
\label{ch:ch08_class_group:8-4-first-computational-labs}

Helpers: \texttt{lib/composition.py} \(\cdot\) \textbf{Appendix C ?C.4}.

\begin{itemize}
\item \textbf{8.A} \texttt{compose\_flywheels} + \texttt{reduce\_composition}.
\item \textbf{8.B} \texttt{composition\_table} / \texttt{closure\_fraction}.
\item \textbf{8.C} \texttt{class\_group\_analogue} order and structure.
\item \textbf{8.D} \texttt{is\_associative\_up\_to\_equivalence} (expect low scores --- OP6 data).
\item \textbf{8.E} Optional: compare composition methods.
\end{itemize}

\begin{lstlisting}[style=qga,language=python]
from lib.composition import compose_flywheels, reduce_composition
# build t1, t2 as in Appendix C ?C.4
# print(reduce_composition(compose_flywheels(t1, t2))["classification"]["type"])
\end{lstlisting}

\bigskip
\noindent\rule{\textwidth}{0.4pt}
\bigskip

\section{Exercises}
\label{ch:ch08_class_group:exercises}

\textbf{8.A (hand).} State Gauss composition for binary quadratic forms in one sentence and explain why it produces a group law on classes.

\textbf{8.B (hand).} Why is associativity (up to equivalence) the key property needed for a class-group analogue?

\textbf{8.C (code).} Complete Labs 8.A--8.B. Report the type of the composed configuration and the composition table's \texttt{closure\_fraction}.

\textbf{8.D (code).} Complete Lab 8.C. What is the size of the class-group analogue on your seed topograph?

\textbf{8.E (code).} Run Lab 8.D. What associativity score do you obtain? Record failures.

\textbf{8.F (Hatcher bridge).} In Hatcher, the class group encodes the failure of unique factorization in quadratic fields. Sketch how a class-group analogue on the Hopf lattice might encode ``failure of unique representation'' by stable flywheels / \(Z\)-map preimages.

\textbf{8.G (Open Problem 6 teaser).} Using \texttt{compose\_flywheels} and \texttt{is\_associative\_up\_to\_equivalence}, test associativity on the current candidate adjacency. Document numerical failures---they \emph{are} OP6 progress.

\textbf{8.H (forward).} Why will the quaternion algebra theory of Chapter 9 be the natural home for a rigorous version of these class groups?

\textbf{8.I (software honesty).} Distinguish: (i) classical class groups (\textbf{Theorem}), (ii) \texttt{class\_group\_analogue} outputs (\textbf{Model} + software), (iii) any claim that Magic Islands \emph{are} class groups of a quaternion order (\textbf{Hypothesis} until proved).

\bigskip
\noindent\rule{\textwidth}{0.4pt}
\bigskip

\section{Code and asset pointers}
\label{ch:ch08_class_group:code-and-asset-pointers}

\begin{lstlisting}[style=qga]
qga/lib/composition.py
qga/lib/flux_topograph.py
qga/lib/hopf_lattice.py
kingdom.core.flux_flywheel   # optional cross-check with Z stability
\end{lstlisting}

\textbf{Figures:} \texttt{scripts/generate\_ch8\_figures.py} \textbf{Open problems:} OP6 (this chapter); depends on OP1--OP3. \textbf{Related portal:} no dedicated composition tab yet---Lattice Simulator / Flux Flywheel remain dynamical counterparts.

\bigskip
\noindent\rule{\textwidth}{0.4pt}
\bigskip

\section{Looking ahead}
\label{ch:ch08_class_group:looking-ahead}

We now have a working \textbf{Model} lift of Gauss composition and the class group to flux configurations on the gauged Hopf lattice---with explicit experimental gaps (closure, associativity). In \textbf{Chapter 9} we move to quaternion algebras and their ideal theory, the rigorous arithmetic home for class groups of orders (Hurwitz, Eichler, ...). \textbf{Chapter 10} returns to observational validation, including whether class-group structure predicts \(350/\pi\) or \(Z\)-map patterns.

With composition and class-group language in place, we are ready for the ideal-theoretic foundation.

\bigskip
\noindent\rule{\textwidth}{0.4pt}
\bigskip


